\documentclass[12pt,a4paper]{article}

\usepackage[utf8]{inputenc}
\usepackage[T1]{fontenc}
\usepackage[spanish]{babel}
\usepackage{geometry}
\usepackage{longtable}
\usepackage{array}
\usepackage{xcolor}
\usepackage{hyperref}
\usepackage{listings}
\usepackage{tikz}

\usetikzlibrary{arrows.meta,positioning,shapes.multipart,fit,backgrounds,babel}

\geometry{margin=2.5cm}

\hypersetup{
    colorlinks=true,
    linkcolor=blue,
    urlcolor=blue
}

\lstdefinestyle{code}{
    basicstyle=\ttfamily\small,
    breaklines=true,
    frame=single,
    columns=fullflexible,
    keepspaces=true,
    showstringspaces=false
}

\title{
    \textbf{Proyecto Final Telegame}\\
    \large Informe tecnico de Teleinformatica
}
\author{
    Estudiante: \underline{\hspace{8cm}}\\
    Codigo: \underline{\hspace{8cm}}\\
    Materia: Teleinformatica
}
\date{Mayo de 2026}

\begin{document}

\maketitle
\newpage

\tableofcontents
\newpage

\section{Introduccion}

Telegame es un sistema cliente-servidor desarrollado en lenguaje C para Linux.
El objetivo del proyecto es implementar un servicio de red donde varios jugadores
puedan conectarse a un servidor central, registrarse, consultar participantes,
entrar a una cola de espera y jugar una partida de tres en raya.

La aplicacion fue disenada para cumplir los requisitos de la materia de
Teleinformatica: usar sockets TCP, mantener un servidor central, permitir varios
clientes dinamicos, administrar el juego desde el servidor y documentar el
protocolo utilizado siguiendo una estructura similar a un RFC.

El juego elegido es tres en raya porque permite demostrar claramente los
conceptos importantes del proyecto: registro de jugadores, turnos, validacion de
comandos, validacion de movimientos, tablero, ganador, empate, abandono y
marcador.

\section{Objetivo general}

Implementar un sistema de juego en red con arquitectura cliente-servidor, donde
un servidor central gestione la conexion de jugadores, el emparejamiento, el
estado del juego y la comunicacion mediante un protocolo de texto sobre TCP.

\section{Objetivos especificos}

\begin{itemize}
    \item Implementar un servidor TCP en C para Linux.
    \item Permitir la conexion de n clientes sin definir un numero maximo fijo en
    la estructura principal de jugadores.
    \item Utilizar \texttt{select(2)} para atender multiples sockets sin usar
    \texttt{fork(2)} ni \texttt{pthread}.
    \item Definir un protocolo de texto llamado \texttt{TELEGAME/1.0}.
    \item Registrar jugadores mediante comandos enviados por el cliente.
    \item Mantener una cola de espera para emparejar jugadores.
    \item Administrar completamente el juego desde el servidor.
    \item Mostrar participantes, turnos, tablero, ganador, empate y marcador.
    \item Documentar comandos, respuestas, automata de estados y estructuras de
    datos.
\end{itemize}

\section{Arquitectura del sistema}

El sistema esta formado por un servidor central y varios clientes. Cada cliente
se conecta al servidor usando TCP. El servidor acepta conexiones, recibe
comandos, procesa cada solicitud y envia respuestas segun el protocolo.

No existe comunicacion directa entre jugadores. Todas las decisiones del juego
son tomadas por el servidor central. Esto permite que el servidor mantenga el
estado oficial del tablero y evite inconsistencias entre clientes.

\subsection{Diagrama general}

\begin{figure}[h]
\centering
\begin{tikzpicture}[
    node distance=2.4cm,
    box/.style={
        draw,
        rounded corners=4pt,
        very thick,
        align=center,
        minimum width=3.2cm,
        minimum height=1.45cm,
        fill=blue!6
    },
    server/.style={
        draw,
        rounded corners=4pt,
        very thick,
        align=center,
        minimum width=4.3cm,
        minimum height=1.8cm,
        fill=green!9
    },
    note/.style={
        draw,
        rounded corners=3pt,
        align=center,
        fill=yellow!15,
        font=\small,
        inner sep=5pt
    },
    arrow/.style={
        {Stealth[length=3mm]}-{Stealth[length=3mm]},
        very thick
    }
]

\node[server] (srv) {
    \textbf{Servidor maestro}\\
    TCP :5000\\
    \texttt{select + juego}
};

\node[box, left=3.4cm of srv] (a) {
    \textbf{Cliente A}\\
    Jugador X\\
    Ramiro
};

\node[box, right=3.4cm of srv] (b) {
    \textbf{Cliente B}\\
    Jugador O\\
    Paulo
};

\node[box, below=2.6cm of srv] (n) {
    \textbf{Cliente N}\\
    Lobby / consulta\\
    Carla
};

\node[note, above=1.1cm of srv] (logic) {
    El servidor central valida\\
    turnos, movimientos y resultados
};

\draw[arrow] (a) -- node[above]{TCP} node[below]{comandos/respuestas} (srv);
\draw[arrow] (srv) -- node[above]{TCP} node[below]{comandos/respuestas} (b);
\draw[arrow] (n) -- node[right]{TCP} (srv);
\draw[very thick, -{Stealth[length=3mm]}] (logic) -- (srv);

\begin{scope}[on background layer]
    \node[
        draw=gray!45,
        rounded corners=8pt,
        fill=gray!5,
        inner sep=12pt,
        fit=(a)(srv)(b)(n)(logic)
    ] {};
\end{scope}

\end{tikzpicture}
\caption{Arquitectura cliente-servidor de Telegame.}
\label{fig:arquitectura-telegame}
\end{figure}

\subsection{Responsabilidades del servidor}

\begin{itemize}
    \item Escuchar en un puerto TCP predefinido.
    \item Aceptar conexiones de clientes.
    \item Registrar nombres de jugadores.
    \item Mantener la lista de jugadores conectados.
    \item Mantener la cola de espera.
    \item Crear partidas 1 vs 1.
    \item Validar comandos y estados.
    \item Validar turnos y movimientos.
    \item Determinar victoria, empate o abandono.
    \item Enviar tablero, turno, resultado y marcador.
    \item Registrar en consola el intercambio de mensajes.
\end{itemize}

\subsection{Responsabilidades del cliente}

\begin{itemize}
    \item Conectarse al servidor central.
    \item Enviar comandos escritos por el jugador.
    \item Recibir mensajes del servidor.
    \item Mostrar el tablero en formato legible.
    \item Mostrar turnos, participantes, resultados y errores.
\end{itemize}

\section{Tecnologias utilizadas}

\begin{longtable}{|p{4cm}|p{10cm}|}
\hline
\textbf{Elemento} & \textbf{Uso en el proyecto} \\
\hline
Lenguaje C & Implementacion del servidor y cliente. \\
\hline
Linux Mint & Sistema operativo objetivo para compilar y ejecutar. \\
\hline
GCC & Compilador usado mediante el \texttt{Makefile}. \\
\hline
Make & Automatiza la compilacion del cliente, servidor y pruebas. \\
\hline
Sockets TCP & Comunicacion confiable cliente-servidor. \\
\hline
\texttt{select(2)} & Multiplexacion de multiples clientes sin procesos ni hilos. \\
\hline
Protocolo de texto & Comandos y respuestas legibles por terminal. \\
\hline
\end{longtable}

\section{Estructura del proyecto}

\begin{lstlisting}[style=code]
.
|-- Makefile
|-- README.md
|-- docs/
|   |-- RFC-TELEGAME.md
|   |-- DEFENSA.md
|   |-- PROMPT_CHATGPT.md
|   `-- INFORME_TELEGAME.tex
|-- scripts/
|   `-- install_telegame_local.sh
|-- src/
|   |-- client.c
|   `-- server.c
`-- tests/
    `-- smoke_test.sh
\end{lstlisting}

\section{Compilacion y ejecucion}

El proyecto incluye un \texttt{Makefile}. Para compilar se ejecuta:

\begin{lstlisting}[style=code]
make
\end{lstlisting}

Esto genera los binarios:

\begin{lstlisting}[style=code]
bin/telegame_server
bin/telegame_client
\end{lstlisting}

Para iniciar el servidor:

\begin{lstlisting}[style=code]
./bin/telegame_server --host 127.0.0.1 --port 5000
\end{lstlisting}

Para iniciar clientes:

\begin{lstlisting}[style=code]
./bin/telegame_client --host 127.0.0.1 --port 5000 --name Ramiro
./bin/telegame_client --host 127.0.0.1 --port 5000 --name Paulo
./bin/telegame_client --host 127.0.0.1 --port 5000 --name Carla
\end{lstlisting}

\section{Protocolo TELEGAME/1.0}

El protocolo \texttt{TELEGAME/1.0} es un protocolo de texto sobre TCP. Cada
mensaje termina con salto de linea \texttt{\textbackslash n}. Esto permite que
el servidor separe comandos completos aun cuando TCP entregue los bytes por
partes.

\subsection{Comandos del cliente}

\begin{longtable}{|p{4cm}|p{10cm}|}
\hline
\textbf{Comando} & \textbf{Descripcion} \\
\hline
\texttt{HELLO <nombre>} & Registra el nombre del jugador. \\
\hline
\texttt{HELP} & Lista comandos disponibles. \\
\hline
\texttt{QUEUE} & Ingresa a la cola de espera. \\
\hline
\texttt{PLAYERS} & Lista participantes conectados y estado. \\
\hline
\texttt{SCORE} & Muestra victorias, empates y derrotas. \\
\hline
\texttt{BOARD} & Solicita tablero y turno actual. \\
\hline
\texttt{MOVE <1-9>} & Realiza una jugada si es el turno del jugador. \\
\hline
\texttt{QUIT} & Cierra la sesion o abandona una partida. \\
\hline
\end{longtable}

\subsection{Respuestas del servidor}

\begin{longtable}{|p{4cm}|p{10cm}|}
\hline
\textbf{Respuesta} & \textbf{Significado} \\
\hline
\texttt{WELCOME} & El servidor acepto la conexion. \\
\hline
\texttt{OK} & Operacion aceptada. \\
\hline
\texttt{ERR} & Operacion rechazada. \\
\hline
\texttt{INFO} & Mensaje informativo. \\
\hline
\texttt{PLAYER} & Participante conectado y estado. \\
\hline
\texttt{SCORE} & Marcador de un jugador. \\
\hline
\texttt{MATCH} & Partida creada y simbolo asignado. \\
\hline
\texttt{BOARD} & Estado del tablero. \\
\hline
\texttt{TURN} & Turno actual. \\
\hline
\texttt{RESULT WIN} & Partida terminada con ganador. \\
\hline
\texttt{RESULT DRAW} & Partida terminada en empate. \\
\hline
\texttt{RESULT ABORT} & Partida terminada por abandono o desconexion. \\
\hline
\texttt{BYE} & Cierre normal de sesion. \\
\hline
\end{longtable}

\section{Automata de estados del jugador}

El jugador pasa por estados definidos segun su interaccion con el servidor:

\begin{lstlisting}[style=code]
[DESCONECTADO]
      |
      | conexion TCP aceptada
      v
[CONECTADO]
      |
      | HELLO <nombre>
      v
[LOBBY] <--------------------------+
      |                            |
      | QUEUE                      | fin de partida
      v                            |
[EN_COLA]                          |
      | servidor encuentra rival   |
      v                            |
[EN_PARTIDA] ----------------------+
      |
      | QUIT / error / desconexion
      v
[DESCONECTADO]
\end{lstlisting}

\section{Estructuras de datos}

\subsection{\texttt{struct Player}}

La estructura \texttt{Player} representa a cada jugador conectado. Esta
estructura esta definida en \texttt{src/server.c}, cerca de la linea 24.

\begin{longtable}{|p{4cm}|p{10cm}|}
\hline
\textbf{Campo} & \textbf{Rol} \\
\hline
\texttt{fd} & Descriptor del socket TCP del cliente. \\
\hline
\texttt{address} & Direccion IP del cliente. \\
\hline
\texttt{port} & Puerto remoto del cliente. \\
\hline
\texttt{name} & Nombre registrado por el jugador. \\
\hline
\texttt{registered} & Indica si el jugador ya envio \texttt{HELLO}. \\
\hline
\texttt{buffer} & Acumula bytes recibidos hasta formar lineas completas. \\
\hline
\texttt{buffer\_len} & Cantidad de bytes usados en el buffer. \\
\hline
\texttt{in\_queue} & Indica si el jugador esta esperando rival. \\
\hline
\texttt{game\_id} & Identificador de la partida activa o cero si no juega. \\
\hline
\texttt{symbol} & Simbolo asignado, \texttt{X} u \texttt{O}. \\
\hline
\texttt{wins, losses, draws} & Marcador efimero del jugador. \\
\hline
\texttt{next} & Enlace en la lista dinamica de jugadores. \\
\hline
\texttt{queue\_next} & Enlace en la cola dinamica de espera. \\
\hline
\end{longtable}

\subsection{\texttt{struct Game}}

La estructura \texttt{Game} representa una partida activa. Esta definida en
\texttt{src/server.c}, cerca de la linea 42.

\begin{longtable}{|p{4cm}|p{10cm}|}
\hline
\textbf{Campo} & \textbf{Rol} \\
\hline
\texttt{id} & Identificador incremental de la partida. \\
\hline
\texttt{x\_player} & Jugador que usa el simbolo \texttt{X}. \\
\hline
\texttt{o\_player} & Jugador que usa el simbolo \texttt{O}. \\
\hline
\texttt{board} & Tablero de nueve posiciones. \\
\hline
\texttt{turn} & Simbolo que debe jugar en el turno actual. \\
\hline
\texttt{next} & Enlace en la lista de partidas activas. \\
\hline
\end{longtable}

\section{Gestion de jugadores y cola}

El servidor mantiene una lista enlazada de jugadores conectados. Esta decision
evita declarar un arreglo fijo de jugadores dentro de la logica principal. Cada
nuevo cliente aceptado se convierte en un nodo de la lista.

La cola de espera tambien se implementa con punteros. Cuando un jugador envia
\texttt{QUEUE}, el servidor verifica que este registrado, que no este en partida
y que no este ya en cola. Luego lo agrega con \texttt{enqueue\_player}. Cuando
hay dos jugadores esperando, \texttt{try\_match\_players} los extrae con
\texttt{dequeue\_player} y crea una partida con \texttt{create\_match}.

\section{Multiplexacion con \texttt{select}}

El servidor usa \texttt{select(2)} en el ciclo principal. En cada iteracion arma
un conjunto de descriptores que incluye:

\begin{itemize}
    \item el socket del servidor, para aceptar clientes nuevos;
    \item todos los sockets de clientes conectados, para leer comandos.
\end{itemize}

Cuando \texttt{select} informa actividad en el socket del servidor, se llama a
\texttt{accept\_player}. Cuando informa actividad en un cliente, se llama a
\texttt{receive\_from\_player}. Esta estrategia permite atender varios clientes
sin crear procesos ni hilos.

\section{Procesamiento de mensajes}

El servidor recibe bytes por TCP mediante \texttt{recv}. Como TCP es un flujo de
bytes, un comando puede llegar dividido o varios comandos pueden llegar juntos.
Por eso cada jugador tiene un \texttt{buffer}. El servidor acumula datos hasta
encontrar un salto de linea. Luego extrae una linea completa y la procesa como
comando.

El flujo es:

\begin{enumerate}
    \item \texttt{receive\_from\_player} recibe bytes.
    \item \texttt{process\_player\_line} limpia saltos de linea.
    \item \texttt{handle\_command} identifica el comando.
    \item Se llama a la funcion concreta: \texttt{handle\_hello},
    \texttt{handle\_queue}, \texttt{handle\_move}, etc.
    \item \texttt{send\_line} elabora y envia la respuesta.
\end{enumerate}

Ademas, el servidor imprime logs de entrada y salida:

\begin{lstlisting}[style=code]
[RX <- Ramiro] MOVE 5
[TX -> Paulo] TURN O Paulo
\end{lstlisting}

Esto ayuda a demostrar en vivo el intercambio de mensajes entre clientes y
servidor.

\section{Logica del juego}

Cada partida tiene dos jugadores. El primero recibe \texttt{X} y el segundo
recibe \texttt{O}. Siempre inicia \texttt{X}. El tablero se representa como una
cadena de nueve caracteres:

\begin{lstlisting}[style=code]
BOARD X.O...X..
\end{lstlisting}

Cada posicion corresponde a una casilla:

\begin{lstlisting}[style=code]
1 | 2 | 3
---------
4 | 5 | 6
---------
7 | 8 | 9
\end{lstlisting}

La funcion \texttt{handle\_move} valida:

\begin{itemize}
    \item que el jugador este en una partida;
    \item que sea su turno;
    \item que el movimiento sea un numero entre 1 y 9;
    \item que la casilla este vacia.
\end{itemize}

Luego actualiza el tablero, avisa a ambos jugadores y verifica si existe ganador
o empate. Si la partida termina, ambos jugadores vuelven al estado \texttt{LOBBY}
y el servidor envia automaticamente participantes y marcador.

\section{Mapa de implementacion}

\begin{longtable}{|p{6cm}|p{8cm}|}
\hline
\textbf{Tema} & \textbf{Ubicacion aproximada} \\
\hline
Estructura \texttt{Player} & \texttt{src/server.c}, linea 24. \\
\hline
Estructura \texttt{Game} & \texttt{src/server.c}, linea 42. \\
\hline
Envio de respuestas & \texttt{send\_line}, linea 80. \\
\hline
Cola de espera & \texttt{enqueue\_player}, linea 159; \texttt{dequeue\_player}, linea 171. \\
\hline
Listado de participantes & \texttt{send\_players\_to}, linea 290. \\
\hline
Marcador & \texttt{send\_score\_to}, linea 280. \\
\hline
Creacion de partida & \texttt{create\_match}, linea 309. \\
\hline
Registro de jugador & \texttt{handle\_hello}, linea 362. \\
\hline
Entrada a cola & \texttt{handle\_queue}, linea 375. \\
\hline
Validacion de jugadas & \texttt{handle\_move}, linea 403. \\
\hline
Interpretacion de comandos & \texttt{handle\_command}, linea 495. \\
\hline
Recepcion de mensajes & \texttt{receive\_from\_player}, linea 560. \\
\hline
Aceptar clientes & \texttt{accept\_player}, linea 602. \\
\hline
Multiplexacion & \texttt{select}, linea 707. \\
\hline
Cliente & \texttt{src/client.c}. \\
\hline
\end{longtable}

\section{Demostracion recomendada}

Para la demostracion se recomienda abrir cuatro terminales: una para el servidor
y tres para clientes.

\subsection{Terminal 1: servidor}

\begin{lstlisting}[style=code]
cd ~/ProyectoTeleinformatica-
make clean
make
./bin/telegame_server --host 127.0.0.1 --port 5000
\end{lstlisting}

\subsection{Terminal 2: cliente Ramiro}

\begin{lstlisting}[style=code]
cd ~/ProyectoTeleinformatica-
./bin/telegame_client --host 127.0.0.1 --port 5000 --name Ramiro
\end{lstlisting}

\subsection{Terminal 3: cliente Paulo}

\begin{lstlisting}[style=code]
cd ~/ProyectoTeleinformatica-
./bin/telegame_client --host 127.0.0.1 --port 5000 --name Paulo
\end{lstlisting}

\subsection{Terminal 4: cliente Carla}

\begin{lstlisting}[style=code]
cd ~/ProyectoTeleinformatica-
./bin/telegame_client --host 127.0.0.1 --port 5000 --name Carla
\end{lstlisting}

Pasos para mostrar:

\begin{enumerate}
    \item En Carla escribir \texttt{PLAYERS} para mostrar tres clientes.
    \item En Ramiro y Paulo escribir \texttt{QUEUE}.
    \item Mostrar que el servidor asigna \texttt{X} y \texttt{O}.
    \item Intentar un movimiento fuera de turno.
    \item Jugar una partida completa.
    \item Mostrar \texttt{RESULT WIN} o \texttt{RESULT DRAW}.
    \item Mostrar que al final ambos jugadores vuelven a \texttt{LOBBY}.
    \item Mostrar el marcador con \texttt{SCORE}.
\end{enumerate}

\section{Pruebas}

El proyecto incluye una prueba automatica basica:

\begin{lstlisting}[style=code]
make test
\end{lstlisting}

La prueba compila el servidor y el cliente, levanta el servidor en
\texttt{127.0.0.1}, conecta dos jugadores, verifica el registro, prueba un
movimiento fuera de turno, ejecuta una partida y valida el ganador y el marcador.

\section{Referencias tecnicas}

Las funciones principales investigadas pertenecen a POSIX sockets:

\begin{itemize}
    \item \texttt{socket(2)}: crea un socket.
    \item \texttt{setsockopt(2)}: configura opciones como \texttt{SO\_REUSEADDR}.
    \item \texttt{bind(2)}: asocia IP y puerto.
    \item \texttt{listen(2)}: pone el socket en modo escucha.
    \item \texttt{accept(2)}: acepta clientes.
    \item \texttt{select(2)}: multiplexa descriptores.
    \item \texttt{recv(2)}: recibe datos.
    \item \texttt{send(2)}: envia datos.
    \item \texttt{inet\_pton(3)} e \texttt{inet\_ntop(3)}: convierten direcciones IP.
\end{itemize}

En Linux Mint se pueden consultar con:

\begin{lstlisting}[style=code]
man 2 socket
man 2 bind
man 2 listen
man 2 accept
man 2 select
man 2 recv
man 2 send
man 3 inet_pton
\end{lstlisting}

\section{Uso de inteligencia artificial}

Durante el desarrollo se utilizo asistencia de IA como apoyo para revisar
requisitos, ordenar la documentacion, comparar alternativas y preparar material
de defensa. La IA no reemplaza la comprension del codigo: el estudiante debe
poder explicar la arquitectura, el protocolo, las estructuras y las funciones
principales.

Para completar esta seccion con informacion personal, se recomienda indicar:

\begin{itemize}
    \item que herramienta de IA se utilizo;
    \item si fue version gratuita o de pago;
    \item en que partes ayudo;
    \item que partes fueron revisadas personalmente;
    \item que aprendizajes quedaron del proyecto.
\end{itemize}

Ejemplo de redaccion adaptable:

\begin{quote}
Use una herramienta de inteligencia artificial como apoyo para organizar ideas,
revisar requisitos y mejorar la documentacion. Los aportes fueron pertinentes
cuando ayudaron a ordenar el protocolo, identificar funciones importantes y
preparar la defensa. Sin embargo, fue necesario revisar el codigo y ejecutar las
pruebas para comprobar que el programa cumpliera los requisitos. El aprendizaje
principal fue entender como un servidor TCP puede atender varios clientes con
\texttt{select}, como se disena un protocolo de texto y como se mantiene el
estado del juego desde un servidor central.
\end{quote}

\section{Conclusiones}

\begin{itemize}
    \item El proyecto cumple con una arquitectura cliente-servidor donde el
    servidor central administra completamente el juego.
    \item El uso de \texttt{select} permite manejar varios clientes sin usar
    procesos ni hilos.
    \item El protocolo de texto facilita la prueba manual, la depuracion y la
    explicacion durante la defensa.
    \item Las estructuras \texttt{Player} y \texttt{Game} separan claramente los
    datos de jugadores y partidas.
    \item La cola de espera permite emparejar jugadores de forma simple y
    dinamica.
    \item El registro y marcador son efimeros, es decir, existen mientras el
    servidor esta ejecutandose.
    \item La demostracion con tres clientes permite mostrar que dos pueden jugar
    mientras otro sigue interactuando con el servidor.
\end{itemize}

\section{Aprendizajes alcanzados}

\begin{itemize}
    \item Uso de sockets TCP en C.
    \item Diferencia entre \texttt{127.0.0.1} para prueba local y
    \texttt{0.0.0.0} para escuchar desde otras computadoras.
    \item Manejo de multiples clientes con \texttt{select}.
    \item Diseno de comandos y respuestas en un protocolo.
    \item Validacion de estados antes de ejecutar acciones.
    \item Uso de listas enlazadas para representar jugadores y cola.
    \item Importancia de documentar el codigo con funciones, estructuras y
    comportamiento observable.
\end{itemize}

\end{document}
